\documentclass[11pt]{article}
\usepackage[a4paper,margin=1in]{geometry}
\usepackage{fontspec}
\usepackage[boldfont=false]{xeCJK}
\setCJKmainfont{FandolSong-Regular.otf}
\usepackage{amsmath}
\usepackage{amssymb}
\usepackage{array}
\usepackage{booktabs}
\usepackage{graphicx}
\usepackage{adjustbox}
\usepackage{enumitem}
\setlist[itemize]{topsep=2pt,partopsep=0pt,parsep=0pt,itemsep=2pt,leftmargin=1.4em}
\setlist[enumerate]{topsep=2pt,partopsep=0pt,parsep=0pt,itemsep=2pt,leftmargin=1.6em}
\usepackage{parskip}
\usepackage{fvextra}
\fvset{breaklines=true,breakanywhere=true,fontsize=\small}
\setlength{\parindent}{0pt}

\begin{document}

\begin{center}
  {\LARGE\bfseries The basic economic problem}\\[0.35em]
  {\large Igcse Economics}
\end{center}
\vspace{0.6em}

\section*{The economic problem}

Economics is about how people use \textbf{resources} 资源 to meet their needs and wants. The main idea is simple. People want a lot, but there are not enough resources to make everything. This is the \textbf{economic problem} 经济问题.

Two facts cause the economic problem:

\begin{itemize}
\item Resources are \textbf{finite} 有限的. There is only a limited amount of them.

\item People's \textbf{wants} 欲望 are unlimited. People always want more and better things.

\end{itemize}

Because wants are unlimited but resources are limited, resources are \textbf{scarce} 稀缺. We call this problem \textbf{scarcity} 稀缺性. Scarcity is the root of all economics.

\subsection*{Goods, services, and choices}

We use resources to make \textbf{goods} 物品 (things you can touch, like food and phones) and \textbf{services} 服务 (jobs done for you, like a haircut or a bus ride).

Because resources are scarce, we cannot have everything. So everyone must make \textbf{choices}. Scarcity forces four groups to choose:

\begin{itemize}
\item \textbf{consumers} 消费者 — people who buy and use goods and services. They choose what to buy.

\item \textbf{workers} 工人 — they choose which job to do.

\item \textbf{producers} 生产者, also called \textbf{firms} 企业 — they make goods and services. They choose what to produce.

\item \textbf{governments} 政府 — they choose what to provide and how to spend public money.

\end{itemize}

\subsection*{Economic goods and free goods}

Most things are \textbf{economic goods} 经济物品. They are scarce, so making or using one means giving something else up. Economic goods have an \textbf{opportunity cost} 机会成本 (explained below).

A few things are \textbf{free goods} 免费物品. They are not scarce — there is enough for everyone and they cost nothing to use. Examples are the air you breathe and sunlight. A free good has no opportunity cost.

Be careful: a "free" sample or a "free" gift is \emph{not} a free good. Resources were used to make it, so it is still an economic good.

\section*{Factors of production}

To make goods and services, firms use four types of resource. We call them the \textbf{factors of production} 生产要素.

\subsection*{The four factors}

\begin{center}
\adjustbox{max width=\linewidth}{%
\begin{tabular}{lll}
\toprule
\textbf{Factor} & \textbf{What it is} & \textbf{Example} \\
\midrule
\textbf{land} 土地 & natural resources & fields, water, oil, forests \\
\textbf{labour} 劳动 & human work & a teacher, a builder, a nurse \\
\textbf{capital} 资本 & man-made things used to produce & machines, tools, factories \\
\textbf{enterprise} 企业家才能 & organising the other three and taking risks & a person who starts a business \\
\bottomrule
\end{tabular}}
\end{center}

The person who provides enterprise is the \textbf{entrepreneur} 企业家. They bring the other factors together and risk their own money to start and run a business.

\subsection*{The reward to each factor}

Each factor earns a \textbf{reward} 报酬 — the income it gets for being used:

\begin{center}
\adjustbox{max width=\linewidth}{%
\begin{tabular}{ll}
\toprule
\textbf{Factor} & \textbf{Reward} \\
\midrule
land & \textbf{rent} 地租 \\
labour & \textbf{wages} 工资 \\
capital & \textbf{interest} 利息 \\
enterprise & \textbf{profit} 利润 \\
\bottomrule
\end{tabular}}
\end{center}

\subsection*{Mobility, quantity and quality}

The \textbf{mobility} 流动性 of a factor means how easily it can move to a different use or place. A worker with many skills can change jobs easily, so labour is mobile. A worker with only one skill, or a factory that cannot be moved, is less mobile.

The quantity and quality of factors can change:

\begin{itemize}
\item \textbf{Quantity} can rise or fall. The amount of labour grows if the birth rate rises or people move into the country to work. Land can be lost to the sea or gained by draining it.

\item \textbf{Quality} can rise. Training and education make labour better. New technology makes capital better. Higher quality means each factor can produce more.

\end{itemize}

\section*{Opportunity cost}

When you choose one thing, you give up the next best choice. That is the real cost of your decision.

\textbf{Opportunity cost} is the next best \textbf{alternative} 替代选择 you give up when you make a choice.

It affects every decision-maker:

\begin{itemize}
\item A \textbf{consumer} who buys a game gives up the book they could have bought instead.

\item A \textbf{worker} who takes one job gives up the pay of the next job they could have done.

\item A \textbf{producer} who makes more cars uses factors that could have made trucks.

\item A \textbf{government} that builds a hospital gives up the school it could have built with the same money.

\end{itemize}

Opportunity cost matters because resources are scarce. Every choice has a cost, even when no money changes hands.

\section*{Production possibility curve (PPC) diagrams}

A \textbf{production possibility curve} 生产可能性曲线 (PPC) shows the largest amounts of two goods a country can make when it uses all its resources fully.

Imagine an economy that makes only two types of good: capital goods and consumer goods. (Capital goods are used to make other goods, like machines. Consumer goods are used up by people, like food and clothes.) We put one good on each axis. The PPC is the line joining every point where the economy makes the most it can.

\subsection*{Reading the diagram}

Picture the two goods on the two axes. The PPC is a curve that bends from one axis down to the other (it usually bows outward, away from the corner):

\begin{itemize}
\item A point \textbf{on} the curve: all resources are used fully. The economy makes the most it can.

\item A point \textbf{inside} the curve: resources are wasted or not all used — for example, there is \textbf{unemployment} 失业. The economy could make more.

\item A point \textbf{outside} the curve: this is impossible now. There are not enough resources to reach it.

\end{itemize}

\subsection*{Movements along the curve}

If the economy already uses all its resources (a point on the curve), making more of one good means making less of the other. You must move factors from one good to the other. This is a \textbf{trade-off} 取舍.

The amount of the second good you give up is the opportunity cost of the extra first good. This is why the PPC is a clear picture of opportunity cost.

\subsection*{Shifts of the curve}

The whole curve can move:

\begin{itemize}
\item It shifts \textbf{outward} (to the right) when the country can make more of \emph{both} goods. This is \textbf{economic growth} 经济增长. Causes: more resources (a bigger workforce, more land discovered) or better-quality resources (training, new technology).

\item It shifts \textbf{inward} (to the left) when the country can make less. This is \textbf{economic decline} 经济衰退. Causes: fewer or worse resources — for example a war, a natural disaster, or using up \textbf{non-renewable resources} 不可再生资源 (resources like oil and coal that cannot be replaced once used).

\end{itemize}

Do not confuse the two. Moving from inside the curve \emph{to} the curve (by using idle resources) is a short-term gain. Moving the whole curve outward is a long-term change.


\end{document}
